\chapter{Requirements}
\label{ch:requirements}

This chapter is dedicated for the requirements of the Zinc compiler.
Each item carries a stable identifier:
\begin{itemize}
    \item \texttt{FR-xx}: Function Requirements.
    \item \texttt{NFR-xx}: Non-Function Requirements.
    \item \texttt{AC-xx}: Acceptance Criteria.
\end{itemize}

\section{Scope and Stakeholders}

Zinc is a compiler that translates \texttt{.zn} source files into native executables through LLVM. The steakeholders are users who write Zinc programs (tipically systems-level code with selective C interoperability) and constributors who extend the compiler itself.

The following are explicitly \emph{out of scope}:

\begin{itemize}
    \item Just-in-time execution and an interactive REPL.
    \item A package manager or remote dependency resolver.
    \item Targets other than what the host LLVM toolchain natively supports.
    \item A formal language specification; the language reference (chapter \ref{ch:requirements} onwards) is the operative description.
\end{itemize}

\section{Functional Requirements}
\label{sec:functional_requirements}

Functional requirements describe the externally observable behaviour of the compiler. The are grouped by pipeline phase.

\subsection{Compiler Pipeline}

\begin{description}
    \item[FR-01 (Source ingestion)] The compiler must accept one or more \texttt{.zn} source files passed on the command line and resolve their \zn{use} imports transitively before any other phase runs.
    \item[FR-02 (Lexical analysis)] The compiler must convert each source file into a stream of \texttt{ZToken} values, preserving the source location for every token (see chapter \ref{ch:lexer}).
    \item[FR-03 (Syntax analysis)] The compiler must build an AST of \texttt{ZNode} values via recursive-descent parsing and must recover from a malformed statement without aborting the run (see chapter \ref{ch:parser}).
    \item[FR-04 (Semantic analysis)] The compiler must perform name resolution, type checking, scope management and capability checking on the AST before the code generation (see chapter \ref{ch:semantics}).
    \item[FR-05 (Code generation)] The compiler must emit LLVM IR for every type-checked module and must produce a native executable by invoking the linker \texttt{lld} on the resulting objects (see chapther \ref{ch:codegen}).
    \item[FR-06 (CLI layer)] The compiler must expose at least the flags required to:emit textual IR, select the output binary path, print compiler logs and provide an help command that shows how to use the CLI.
\end{description}

\subsection{Language Support}

\begin{description}
    \item[FR-07 (Primitive types)] The compiler must accept the primitive types (see chapter \ref{ch:language_reference}).
    \item[FR-08 (Compund types)] The compiler must accept pointer types, fixed arrays and slices, tuples and function types.
    \item[FR-09 (User-defined types)] The compiler must aceept \zn{struct} and \zn{enum} declarations, including struct embedding and enum payloads, and must rejects infinite-size types.
    \item[FR-10 (Function and methods)] The compiler must accept free functions, receiver methods (\zn{T.method} and \zn{T::method}).
    \item[FR-11 (Control flow)] The compiler must accept \zn{if}/\zn{else} C-style and while-style loops.
    \item[FR-12 (C interoperability)] The compiler must accept \zn{foreign} declarations and must emit calls compatile with the target C ABI so that Zinc code can call into the C standard library.
\end{description}

\subsection{Diagnostics}

\begin{description}
    \item[FR-15 (Error location)] Every reported error or warning must point at a concrete source location (file, row, col).
    \item[FR-16 (Error continuation)] On encountering an error in any phae, the compiler must continue alaysing the remaining input without propagating the error in the entire source code.
    \item[FR-17 (Unused symbols warnings)] The compiler must emit warnings about unused variables/functions and types that are declared but never referenced, unless suppressed by a CLI flag.
\end{description}

\section{Use Cases}
\label{sec:use_cases}

\begin{description}
    \item[UC-01 (Compile to native executable)] Actor: developer. Main flow: runs \texttt{zinc file.zn -o out}; the compiler executes the full pipeline and writes a native binary. Error flow: any phase failure emits diagnostics with source location (FR-15) and the compiler exits non-zero without producing the binary.
    \item[UC-02 (Emit LLVM IR)] Actor: developer. Main flow: runs \texttt{zinc file.zn --emit-llvm -o file.ll}; the front-end runs and writes textual IR. Error flow: invalid input produces diagnostics and a non-zero exit; no \texttt{.ll} file is written.
\end{description}

\section{Non-Functional Requirements}

Non-functional requirements describe qualities the implementation must hold regardless of which functional requirements are exercised.

\begin{description}
    \item[NFR-01 (Implementation language)] The compiler must be written in C, with the linker driver in C++, and must build with any C11-capable toolchain on the supported platforms.
    \item[NFR-02 (Portability)] The compiler must build and run on macOS and Linux on \texttt{x86\_64} and \texttt{arm64}, using the platform's LLVM installation.
    \item[NFR-03 (Determinism)] Given the same inputs, the same CLI flags and the same LLVM version, the compiler must produce the same output.
    \item[NFR-04 (Memory)] All per-compilation allocations must go through the arena allocator so that a single compilation unit can be freed in one operation; \texttt{malloc} is disallowed.
    \item[NFR-05 (Maintainability)] The compiler must keep each pipeline phase in its own translation unit and must shate state through a signel \texttt{ZState} value rather than through globals.
    \item[NFR-06 (Documentation)] The behaviour described in this document must match the implementation; when the two diverge, the document is updated in the same change set that introduces the divergence.
\end{description}

\section{Acceptance Criteria}
\label{sec:acceptance_criteria}

Acceptance criteria are written so that each one can be checked by an automated test or a short manual procedure. They are referenced from the test suite via the test file name where applicable.

\begin{description}
    \item[AC-01 (e2e simple program)] Compiling the program in section \emph{First Program} of chapter \ref{ch:requirements} with the default flags must produce an executable that prints \texttt{Hello world} followed by a newline and exits with status \texttt{0}.
    \item[AC-02 (Pass test)] Everi fule under \texttt{tests/pass} must compile without errors and , when run, must exit with status \texttt{0}.
    \item[AC-03 (Fail test)] Every file under \texttt{tests/fail} must be rejected by the compiler with at least one diagnostic that points at a valid source location. The compiler must not crash on any of these inputs.
    \item[AC-04 (Performance)] A synthesised 10,000 line input compiled with \texttt{--emit-llvm} must complete the front-end in under one second on the reference hardware described in NFR-05.
\end{description}

\begin{note}
Acceptance criteria are intentionally weaker than a full specification because they state what a pasing implementation must demonstrate, not the full set of behaviours that a finished compiler must exhibit.
\end{note}
